\begin{table}[!htbp]
\centering
\small
\caption{Model Comparison: Log-Likelihood and Information Criteria}
\label{tab:model_comp}
\begin{threeparttable}
    \begin{tabular}{@{}llllll@{}}
        \toprule
        Model & Log-lik. & AIC & BIC & $p_{\mathrm{eff}}$ & Notes \\
        \midrule
        OLS (FR only) & -109.3 & 224.7 & 229.9 & 3.0 & Gaussian OLS baseline \\
        SSM (fixed $\beta$) & -108.6 & 229.2 & 239.7 & 6.0 & OLS with controls \\
        TVP-SSM (EPU+GPR) & -107.1 & 316.1 & 406.0 & 51.0 & Posterior-mean TVP \\
        \bottomrule
    \end{tabular}
    \begin{tablenotes}[flushleft]
        \footnotesize
        \item \textit{Notes:} Log-lik.\ is total-sample log-likelihood evaluated at point estimates.
        \item For TVP-SSM, log-lik.\ uses posterior-mean parameters and the $\beta_t$ mean path.
        \item Controls are standardized as in the TVP-SSM; $FE$/$FR$ are in percentage-point units.
        \item AIC/BIC are computed from log-lik.\ and $p_{\mathrm{eff}}$; $p_{\mathrm{eff}}$ counts time-varying $\beta_t$ for TVP.
        \item TVP-SSM parameters use posterior means from Table 7 ($R_{\text{var}}=9.721$, $Q_{\text{var}}=1.003$).
    \end{tablenotes}
\end{threeparttable}
\end{table}
